\section{Universal Formal Group Laws}\label{sec:chromatichomotopy-formalgrouplaw}
% Write your notes here. Do not add a documentclass or document environment.
Recall the definition of a formal group law.

\begin{definition}
    A \textbf{formal group law} over a ring $R$ is a formal power series $F(x,y)$ such that 
    \begin{itemize}
        \item $F(x,0)=F(0,x)=x$
        \item $F(x,y)=F(y,x)$
        \item $F(x,F(y,z))=F(F(x,y),z)$
    \end{itemize}
For $R$ graded, we require $F(x,y)$ to be homogeneous of degree $2$.
\end{definition}
\begin{remark}
    Given an ungraded formal group law $F$ over $R$ determines a continuous ring map $R[[t]]\to R[[x,y]]$, because, $F(x,y)$ has no constant term, so $F(x,y)^n\to 0$. Therefore, there exists a map
     $\mathrm{Spf}(R[[x]])\times \mathrm{Spf}(R[[y]])\to \mathrm{Spf}(R[[t]])$, which is by the above definition, is a group object in the category of affine formal schemes. More on this later.
\end{remark}

\begin{definition}
    Let $F$ be a formal group law over $R$ and $G$ be a formal group law over $S$, a \textbf{morphism of formal group law} $\varphi: F/R\to G/S$ contains a ring map $\psi: R\to S$, and a formal power series $f(t)\in S[[t]]$ such that 
    \begin{align*}
        f(\psi^*F(x,y))=G(f(x),f(y))
    \end{align*}
    where $\psi^*F$ is the formal group law over $S$ by taking the coefficient of $F$ to $S$ along $\psi$.  

    An \textbf{isomorphism of formal group law} is a morphism of formal group law with an inverse, and a \textbf{strict isomorphism of formal group law} is an isomorphism $\varphi=(\psi, f): F/R\to G/S$, such that $f'(0)=1$.
\end{definition}

Given formal group law $F$, we denote $x+_F y:=F(x,y)$.

\begin{definition}
    The functor $\mathrm{fgl}: \Ring\to \Sets$ is such that 
    \[
    \mathrm{fgl}(R)=\{F\in R[[x,y]]: F \text{ is a formal group law over $R$}\}
    \]and for $psi: R\to S$, $\mathrm{fgl}(\psi):\mathrm{fgl}(R)\to \mathrm{fgl}(S)$ takes $F$ to $\psi^*F$.

    Given two formal group laws $F/R$ and $G/S$, the functor $\mathrm{iso}(F,G): \Ring\to \Sets$ is such that 
    \[
    \mathrm{iso}(F,G)(T):=\{(\psi, \phi, f): \psi: R\to T, \phi: S\to T, f\in T[[t]], \text{ $f$ is an isomorphism $f: \psi^*F\to \phi^*G$ }\}
    \]and for $\mu: T\to T'$, $\mathrm{iso}(F,G)(\mu)$ takes $(\psi, \phi, f)$ to $(\mu\circ \psi, \mu\circ \phi, \mu^*f)$.

    Lastly, the functor $\mathrm{strict}(F,G)$ is the subfunctor of $\mathrm{iso}(F,G)$ where $f$ is a strict isomorphism (this does define a subfunctor as unit is preserved under ring map).
\end{definition}

Of course, every definitions here can be generalized for graded ring and graded formalg group law.
\begin{theorem}
\cite{ChrHopkins}
The functor $\mathrm{fgl}$ is corepresentable by $L$.
\end{theorem}

Recall that given a complex orientation on $E$, there exists a graded formal group law over $E^*$ determined by the complex orientation.
